\documentclass[10pt,a4paper]{amsart}
\usepackage[margin=19mm]{geometry}
\usepackage{amsmath,amssymb,mathtools}
\usepackage[scr=rsfs,cal=euler]{mathalfa}
\usepackage{xfrac}
\usepackage[unicode,hidelinks,bookmarks=false]{hyperref}
\newcommand{\ie}{i.e.\ }
\newcommand{\cf}{cf.\ }
\newcommand{\resp}{respectively\ }
\newcommand{\Subsection}{\subsection}
\providecommand{\nobreakdash}{\nobreak\hbox{-}}
\setlength{\emergencystretch}{2em}
\multlinegap0pt
\usepackage{amssymb}
\newtheorem{theorem}{Theorem}
\title{A Variance Representation Formula for H\"older's Inequality}
\author{Li-Chang Hung}
\date{Source: arXiv:2607.23255v1 (CC BY 4.0)}
\begin{document}
\maketitle

\noindent\emph{Source and licence.} A Variance Representation Formula for H\"older's Inequality. Version arXiv:2607.23255v1 (CC BY 4.0), distributed by arXiv under the Creative Commons Attribution 4.0 International licence, http://creativecommons.org/licenses/by/4.0/, as recorded in the arXiv licence metadata for this exact version and shown on the abstract page https://arxiv.org/abs/2607.23255v1. Full licence notice: \texttt{licenses/arxiv-2607.23255-CC-BY-4.0.txt}.\par\smallskip

\noindent\textit{Editorial scope.} Counted unit: the article's theorem theorem (source0.tex lines 432--525). The article's complete original proof of it is delivered with it (source0.tex lines 528--929, one \texttt{proof} environment of 402 lines). No mathematics is written, repaired or re-derived: the statement and the proof are the article's own lines, in the article's own notation, and no retained line is substituted or re-flowed.  No further source-local context is needed: every cross-reference of every retained line resolves inside the two delivered spans.  Nothing was omitted from any retained span, so the package carries no omission record.  No retained line contains a citation, so the wrapper carries no bibliography and no bibliography entry.  No retained line contains a figure environment, an image-inclusion command or drawing code, so no image is delivered and the package declares no asset dependency.  Editorial formatting only, in addition: an AMS \texttt{amsart} preamble that restates the article's own declarations and macros used by the retained lines, namely the environments \texttt{theorem} on separate counters, together with the standard packages those lines need.  The retained environments print in this document as Theorem 1 in order of appearance, whereas the article prints the same items with the numbering of its own full document.  The complete native source of the article as distributed in the arXiv e-print for this exact version is bundled unchanged as \texttt{sources/206013/source0.tex}.
\begin{theorem}[Exact Variance Representation of Hölder Deficit]
Let

\[
1<p,q<\infty,
\qquad
\frac1p+\frac1q=1,
\]

and let

\[
f,g\geq0
\]

satisfy

\[
0<
\int f^p dx<\infty ,
\qquad
0<
\int g^q dx<\infty .
\]

Define

\[
F=f^p,
\qquad
G=g^q,
\]

and

\[
Z(\theta)
=
\int
F^{1-\theta}G^\theta dx .
\]

Let

\[
d\mu_\theta
=
\frac{
F^{1-\theta}G^\theta
}{
Z(\theta)
}
dx .
\]

Then

\[
\boxed{
\log
\frac{
\|f\|_{L^p}
\|g\|_{L^q}
}{
\int fg dx
}
=
\int_0^1
K_{1/q}(s)
\operatorname{Var}_{\mu_s}
\left(
\log G-\log F
\right)
ds ,
}
\]

where

\[
K_{1/q}(s)
=
\begin{cases}
\dfrac{s}{p},
&
0\leq s\leq\dfrac1q ,
\\[2ex]
\dfrac{1-s}{q},
&
\dfrac1q\leq s\leq1 .
\end{cases}
\]

\end{theorem}

\begin{proof}

Define

\[
\Phi(x,\theta)
=
(1-\theta)\log F(x)
+
\theta\log G(x).
\]

Then

\[
Z(\theta)
=
\int e^{\Phi(x,\theta)}dx .
\]

Since $\Phi$ is linear in $\theta$, we have

\[
\Phi_{\theta\theta}=0 .
\]


We first compute the first derivative of $\log Z$.

Differentiating under the integral sign gives

\[
Z'(\theta)
=
\int
\Phi_\theta e^\Phi dx .
\]

Therefore,

\[
\frac{Z'(\theta)}{Z(\theta)}
=
\frac{
\int
\Phi_\theta e^\Phi dx
}{
\int e^\Phi dx
}.
\]

Using the probability measure $\mu_\theta$, this becomes

\[
(\log Z)'(\theta)
=
E_{\mu_\theta}
(\Phi_\theta).
\]

Because

\[
\Phi_\theta
=
\log G-\log F ,
\]

we obtain

\[
(\log Z)'(\theta)
=
E_{\mu_\theta}
(
\log G-\log F
).
\]


We next differentiate once more.

For a general function $H(x,\theta)$, we have

\[
\frac d{d\theta}
E_{\mu_\theta}(H)
=
E_{\mu_\theta}(H_\theta)
+
\operatorname{Cov}_{\mu_\theta}
(H,\Phi_\theta).
\]

Indeed,

\[
E_{\mu_\theta}(H)
=
\frac{
\int H e^\Phi dx
}{
Z(\theta)
}.
\]

Differentiating,

\[
\begin{aligned}
\frac d{d\theta}E_{\mu_\theta}(H)
&=
\frac{
\int(H_\theta+H\Phi_\theta)e^\Phi dx
}{Z}
\\
&\quad
-
\frac{Z'}{Z}
\frac{\int He^\Phi dx}{Z}.
\end{aligned}
\]

Hence

\[
\frac d{d\theta}E_{\mu_\theta}(H)
=
E(H_\theta)
+
E(H\Phi_\theta)
-
E(H)E(\Phi_\theta),
\]

which gives

\[
\frac d{d\theta}E_{\mu_\theta}(H)
=
E(H_\theta)
+
\operatorname{Cov}(H,\Phi_\theta).
\]

Applying this formula with

\[
H=\Phi_\theta ,
\]

we obtain

\[
(\log Z)''
=
E(\Phi_{\theta\theta})
+
\operatorname{Cov}
(\Phi_\theta,\Phi_\theta).
\]

Since

\[
\Phi_{\theta\theta}=0 ,
\]

it follows that

\[
\boxed{
(\log Z)''
=
\operatorname{Var}_{\mu_\theta}
(
\log G-\log F
).
}
\]

This is the fundamental variance identity.


We now convert this local identity into a global representation.

For a twice differentiable function $h$ on $[0,1]$,
the Green representation formula gives

\[
h(\theta)
=
(1-\theta)h(0)
+
\theta h(1)
-
\int_0^1
K_\theta(s)h''(s)ds ,
\]

where

\[
K_\theta(s)
=
\begin{cases}
(1-\theta)s,
&
s<\theta ,
\\[1ex]
\theta(1-s),
&
s>\theta .
\end{cases}
\]


Indeed, let

\[
r(\theta)
=
h(\theta)
-
(1-\theta)h(0)
-
\theta h(1).
\]

Then

\[
r(0)=r(1)=0 .
\]

The Green function for the operator
$-\frac{d^2}{d\theta^2}$ on $[0,1]$
with zero boundary values is

\[
G_\theta(s)
=
\begin{cases}
(1-\theta)s,
&
s<\theta ,
\\[1ex]
\theta(1-s),
&
s>\theta .
\end{cases}
\]

Hence

\[
r(\theta)
=
-\int_0^1
G_\theta(s)h''(s)ds ,
\]

which proves the formula.


We apply this representation to

\[
h(\theta)=\log Z(\theta).
\]

Taking

\[
\theta=\frac1q ,
\]

we obtain

\[
\log Z(1/q)
=
\frac1p\log Z(0)
+
\frac1q\log Z(1)
-
\int_0^1
K_{1/q}(s)
(\log Z)''(s)ds .
\]


The endpoint values are

\[
Z(0)
=
\int f^p dx ,
\]

and therefore

\[
\log Z(0)
=
p\log\|f\|_p .
\]

Similarly,

\[
\log Z(1)
=
q\log\|g\|_q .
\]

Moreover,

\[
Z(1/q)
=
\int
(f^p)^{1/p}
(g^q)^{1/q}dx ,
\]

and hence

\[
Z(1/q)
=
\int fg dx .
\]


Therefore,

\[
\begin{aligned}
\log\int fg dx
&=
\frac1p
p\log\|f\|_p
+
\frac1q
q\log\|g\|_q
\\
&\quad
-
\int_0^1
K_{1/q}(s)
(\log Z)''(s)ds .
\end{aligned}
\]

Thus

\[
\log\int fg dx
=
\log\|f\|_p
+
\log\|g\|_q
-
\int_0^1
K_{1/q}(s)
(\log Z)''(s)ds .
\]

Using the variance identity,

\[
(\log Z)''(s)
=
\operatorname{Var}_{\mu_s}
(
\log G-\log F
),
\]

we conclude that

\[
\log
\frac{
\|f\|_p\|g\|_q
}{
\int fg dx
}
=
\int_0^1
K_{1/q}(s)
\operatorname{Var}_{\mu_s}
(
\log G-\log F
)
ds .
\]

The proof is complete.

\end{proof}

\end{document}
